\chaptoc{Introduction}

\sectoc{Problem Description}

The goal of an $n \times n$ Sudoku puzzle is to fill each cell with a number $1$ to $n$,
so that all rows, columns, and regions have the same number exactly once.
Variant of a Sudoku adds some additional rules.

\sectoc{State of the Art}

There are multiple ways to solve a Sudoku puzzle.

\subsection*{Human-style logical strategies}

Humans solve Sudoku using easy-to-explain strategies.
The goal of a strategy is to find a cell that cannot contain some numbers and explain why.
These strategies might not always be easy to find,
because they sometimes require to go through many possible combinations to find them.
However once a strategy is found, it should be easy explain why certain cells cannot contain certain numbers.
Humans can often use intuition to speed up solving,
for example by looking at places with more numbers or searching around numbers that have been recently filled in.

\subsection*{Backtracking Algorithm}

A recursive trial-and-error algorithm that solves an $n \times n$ Sudoku puzzle.
For each empty cell, the algorithm attempts to place every number from $1$ to $n$.
If a placement is valid, the algorithm recurses to the next empty cell.
If a solution is found, it is added to the list of solutions.
If no valid number exists for a cell, the algorithm backtracks to the previous cell and tries the next candidate.

The algorithm can be significantly sped up by applying basic solving strategies Chapter~\ref{sec:basic-strategies},
which eliminate impossible candidates from cells before any trial placement is made.

\sectoc{Goals of the Thesis}

The main goal of this thesis is to develop software that helps people solve Sudoku variants
by applying strategies from simplest to most complex.
A user can configure a Sudoku task in the UI\@.
The software then provides step-by-step explanation of how the puzzle can be solved.
The software is programmed in TypeScript, HTML and CSS and can be opened in an internet browser.

\sectoc{Thesis Structure}

Chapter~\ref{ch:chap-1} explains rules of Sudoku and the variants implemented by the software.
Chapter~\ref{ch:chap-2} explains solving strategies implemented by the software.
These two chapters are meant to be used as part of User documentation,
to explain Sudoku rules, variants, and solving strategies to people who never solved a Sudoku before
and only want to know how to use the software.

Chapters~\ref{ch:3}~and~\ref{ch:4} explain Sudoku rules and solving strategies formally.
They are meant to explain things that are not possible to explain without using mathematical text.
They also set up theory which is used in Chapter~\ref{chap:5}.

Chapter~\ref{chap:5} talks about Perfect matching strategies,
which are strategies that can be explained by converting a part of Sudoku state to a bipartite graph,
and then finding edges that are not part of a perfect matching.
The main objective of this chapter is to prove Algorithm~\ref{alg:algorithm},
which can find the smallest application of Perfect matching strategies in polynomial time.
